\chapter{A-tól Z-ig: egy minimális webalkalmazás}

Az előző fejezetekben felépítettük az RWLang model- és capability-szemléletét, majd lefordítottuk és telepítettük a rendszert. A teljes szerverkonfigurációt még nem kell ismernünk: ehhez a példához csak a minimális development beállításokat használjuk. Mielőtt a nyelv egyes részeit külön-külön részleteznénk, készítsünk el egy teljes, de szándékosan kicsi alkalmazást. A cél nem egy kész webshop, hanem annak megmutatása, hogyan áll össze egy RWLang webalkalmazás a forrásfájltól az adatbázison és a HTTP kérésen át a production telepítésig.

A példa egy mini terméklista. A kezdőoldal kilistázza az adatbázisban lévő termékeket, és ugyanott egy formmal új terméket is fel lehet venni. Ezzel egyetlen példában megjelenik a model, typed SQL, GET page, POST action, form validation, CSRF, transaction, migration, konfiguráció és a reverse proxy mögötti futtatás.

A repositoryban a kész példa itt található:

\begin{lstlisting}
examples/a-z-products/
  main.rw
  migrations/0001_create_products.sql
  server.toml.example
  README.md
\end{lstlisting}

\section{A könyvtár létrehozása}

Fejlesztéshez legyen egy külön alkalmazáskönyvtárunk:

\begin{lstlisting}
mkdir -p my-products/migrations
cd my-products
\end{lstlisting}

Ebben a példában minden RWLang kód egyetlen \code{main.rw} fájlban marad. Nagyobb alkalmazásnál később \code{mod} fájlokra bontjuk.

\section{Az adatbázis-séma}

A \code{migrations/0001_create_products.sql} tartalma SQLite esetén:

\begin{lstlisting}[language=SQL]
CREATE TABLE products (
    id INTEGER PRIMARY KEY AUTOINCREMENT,
    name TEXT NOT NULL,
    price INTEGER NOT NULL CHECK (price >= 0)
);
\end{lstlisting}

Ez a migration hozza létre azt a táblát, amelyet az alkalmazás typed SQL lekérdezései használnak. Productionban a migráció külön deploy lépés; a server nem módosítja automatikusan a sémát induláskor.

\section{A model és a lekérdezések}

A \code{main.rw} elején a modellközpontú fejezetben megismert módon először a typed üzleti adatalakot írjuk le:

\begin{lstlisting}[language=RWLang]
model Product {
    id: Int
    name: String
    price: Int
}
\end{lstlisting}

A lista lekérdezése:

\begin{lstlisting}[language=RWLang]
query fn listProducts(db: Db) -> Result<List<Product>, DbError> sql {
    SELECT id, name, price
    FROM products
    ORDER BY id DESC
    LIMIT 50
}
\end{lstlisting}

A \code{List<Product>} azt mondja ki, hogy a query nulla vagy több \code{Product} sort adhat vissza. A \code{SELECT} oszlopai pontosan a model mezőinek nevét és sorrendjét követik.

Az INSERT transaction capabilityt kap:

\begin{lstlisting}[language=RWLang]
query fn createProduct(
    tx: Transaction,
    name: String,
    price: Int
) -> Result<Void, DbError> sql {
    INSERT INTO products(name, price)
    VALUES (:name, :price)
}
\end{lstlisting}

A \code{:name} és \code{:price} bind paraméter. Nem szövegkonkatenációval építünk SQL-t felhasználói inputból.

\section{A GET oldal}

A kezdőoldal lekéri a termékeket, megjelenít egy formot és kilistázza az eredményt:

\begin{lstlisting}[language=RWLang]
page fn home(ctx: PageContext, db: Db) -> Result<Html, PageError> {
    let products = listProducts(db)?;

    return Ok(html {
        <main>
            <h1>Mini termeklista</h1>

            <form method="post" @action(create)>
                <input type="hidden" name="_csrf" value="{{ csrfToken }}">
                <label>
                    Nev
                    <input name="name" maxlength="100">
                </label>
                <label>
                    Ar
                    <input name="price">
                </label>
                <button type="submit">Hozzaadas</button>
            </form>

            <ul>
                @for product in products {
                    <li>
                        #{{ product.id }} - {{ product.name }} - {{ product.price }}
                    </li>
                }
            </ul>
        </main>
    });
}
\end{lstlisting}

A dinamikus szövegek HTML-escape-elve kerülnek a válaszba. A form action nem kézzel írt URL, hanem az \code{@action(create)} helperből származik. A rejtett \code{_csrf} mező a POST kérés CSRF-védelméhez kell.

\section{A POST action}

A form sikeres beküldése transactionben ír az adatbázisba, majd redirectel vissza a kezdőoldalra:

\begin{lstlisting}[language=RWLang]
action fn create(
    ctx: ActionContext,
    db: Db,
    name: String,
    price: Int
) -> Result<Redirect, PageError> {
    transaction db {
        createProduct(tx, name, price)?;
    }
    return Ok(redirect("/"));
}
\end{lstlisting}

Ez a Post/Redirect/Get minta. Ha a felhasználó frissíti a visszakapott oldalt, a böngésző nem küldi be újra automatikusan az INSERT-et.

\section{A route-ok és a validáció}

A fájl végén két route köti össze a HTTP felületet a handlerekkel:

\begin{lstlisting}[language=RWLang]
route home GET "/" => home;

route create POST "/products"
    form name<String> price<Int>
    validate name length 1 100 price range 0 100000000
    => create;
\end{lstlisting}

A POST route már a handler előtt typed inputtá alakítja a két mezőt. Üres vagy 100 karakternél hosszabb név, negatív ár, túl nagy ár vagy nem egész szám esetén a kérés nem jut el érvényes üzleti műveletként az INSERT-ig.

\section{A teljes alkalmazás egy helyen}

Az előző részek együtt alkotják a teljes alkalmazást: \code{model} + két \code{query fn} + \code{page fn} + \code{action fn} + két \code{route}. Nem ismételjük meg még egyszer ugyanazt a hosszú listinget; a repositoryban a pontos, egyfájlos változat közvetlenül megnyitható:

\begin{lstlisting}
examples/a-z-products/main.rw
\end{lstlisting}

Ez szándékos könyvszerkesztési döntés is: a következő fejezetekben az egyes nyelvi elemeket már részletesen tárgyaljuk, itt viszont az alkalmazás teljes útvonalán akarunk végigmenni.

\section{SQLite kapcsolat developmenthez}

A database URL-t ne írjuk a forrásba. Készítsünk külön fájlt, abszolút SQLite pathszal:

\begin{lstlisting}
printf '%s\n' 'sqlite:///tmp/rwlang-a-z-products.db' > dev-db-url
chmod 600 dev-db-url
\end{lstlisting}

A fájl itt nem valódi titkot tartalmaz, de ugyanazt a \code{*_file} mintát használjuk, mint productionben a jelszavas database URL esetén.

\section{Forrásellenőrzés és migration}

Ha a release binárisokat már telepítettük például \code{/usr/local/bin} alá:

\begin{lstlisting}
/usr/local/bin/rwlang-cli check main.rw

/usr/local/bin/rwlang-cli migrate verify \
  --dir migrations \
  --db-url-file "$PWD/dev-db-url"

/usr/local/bin/rwlang-cli migrate apply \
  --dir migrations \
  --db-url-file "$PWD/dev-db-url"
\end{lstlisting}

A \code{check} a nyelvi és statikus szerződést ellenőrzi. A migration CLI külön folyamat; ettől még az alkalmazásszerver nem indult el.

\section{Minimális development konfiguráció}

A \code{server.toml} legyen például:

\begin{lstlisting}
[server]
app = "/ABS/PATH/my-products/main.rw"
listen = "127.0.0.1:8080"
insecure_dev_cookies = true

[database]
url_file = "/ABS/PATH/my-products/dev-db-url"
\end{lstlisting}

A két \code{/ABS/PATH} helyére valódi abszolút útvonal kell. Az \code{insecure_dev_cookies} csak loopback developmenthez elfogadható; productionben ne kapcsold be.

Először ellenőrizzük a konfigurációt:

\begin{lstlisting}
/usr/local/bin/rwlang-server --config "$PWD/server.toml" --check-config
\end{lstlisting}

Majd indulhat a szerver:

\begin{lstlisting}
/usr/local/bin/rwlang-server --config "$PWD/server.toml"
\end{lstlisting}

\section{Az első HTTP kérés}

Másik terminálból:

\begin{lstlisting}
curl -i http://127.0.0.1:8080/
\end{lstlisting}

Böngészőben nyisd meg ugyanazt az URL-t. A form beküldése után a folyamat:

\begin{enumerate}
  \item a böngésző POST-ot küld a \code{/products} route-ra;
  \item a runtime ellenőrzi a request schema-t, a typed mezőket, a validációt és a CSRF tokent;
  \item a \code{create} action transactiont nyit;
  \item a \code{createProduct} bind paraméterekkel INSERT-et futtat;
  \item siker esetén commit történik;
  \item a handler redirectet ad a \code{/} URL-re;
  \item az új GET már az adatbázisból visszaolvasott sort rendereli.
\end{enumerate}

Ez az egyszerű lánc a későbbi wiki és CMS alkalmazásban is ugyanaz marad; csak több model, query, route, jogosultság és komponens kerül köré.

\section{Mit történik hibás inputnál?}

Próbálj negatív árat vagy nem numerikus értéket megadni. A fontos elv az, hogy a böngésző HTML attribútumai csak felhasználói kényelmet adnak; a valódi szerződést a szerveroldali route schema és validation adja. Az SQL query csak már typed értéket kap.

Érdemes még a következőket kipróbálni:

\begin{itemize}
  \item küldj hiányzó \code{name} mezőt;
  \item küldj ismeretlen extra form mezőt;
  \item módosítsd a migrationt alkalmazás után, majd futtasd újra a \code{migrate verify} parancsot;
  \item állítsd le az adatbázis elérését, és figyeld meg a readiness illetve a szerverlog viselkedését.
\end{itemize}

\section{Developmentből production felé}

A példa alkalmazáskódja ugyanaz marad release környezetben is; nem kell újraírni csak azért, mert a szervert már telepített binárissal futtatjuk. A \\ref{ch:telepites}. fejezetben bemutatott build és telepítés után a szerver tipikusan így indul:

\begin{lstlisting}
/usr/local/bin/rwlang-server \\
  --config /usr/local/etc/rwlang/server.toml
\end{lstlisting}

Az alkalmazás entrypointja lehet például \code{/srv/rwlang/current/main.rw}. Apache vagy Nginx esetén a publikus TLS kapcsolatot a proxy terminálja, az RWLang pedig loopbacken hallgat. A teljes Apache/Nginx mintát nem ismételjük meg itt: a telepítési fejezet mutatja be a trust boundaryt, a könyv végi függelék pedig copy/paste-közeli production konfigurációt ad.

\section{Mit tanultunk ebből az egy példából?}

Egy mindössze néhány tucat soros alkalmazásban már láttuk az RWLang teljes alap útvonalát:

\begin{lstlisting}
HTTP GET
  -> route
  -> page
  -> typed query
  -> database
  -> escaped HTML

HTTP POST
  -> route schema + validation + CSRF
  -> action
  -> transaction
  -> typed SQL bind
  -> commit
  -> redirect
  -> GET
\end{lstlisting}

A következő fejezetek ezt a képet bontják szét. Amikor egy-egy nyelvi vagy runtime részlet technikainak tűnik, érdemes visszagondolni erre a két kérésre: \emph{hol lép be az adat a rendszerbe, és milyen typed/capability határokon halad át, mielőtt tartós állapot vagy HTTP válasz lesz belőle?}
